%% ============================================================================
\section{Redemption Concentration}
\label{app:hhi}

We measure redemption concentration on-chain (Section~\ref{sec:setting}) to ask whether a
few large arbitrageurs caused the run. We attribute USDC burns to the external deposits that
funded them by first-in-first-out matching, so every counted dollar corresponds to a burn.
This attribution covers $89.8\%$ of the \$8.15B gross burned. The remaining $10.2\%$
is the pre-window remainder, burns whose funding deposit arrived before the frozen
window opens. Over the $776$ funding addresses,
Table~\ref{tab:hhi} reports both the raw index, $\mathrm{HHI}=\sum_i s_i^2$ over
redemption-volume shares $s_i$, and a size-corrected normalized variant,
$(\mathrm{HHI}-1/N)/(1-1/N)$, which removes the floor a finite $N$ imposes. At $N=776$, the two
differ by under $0.001$ ($0.2558$ raw vs.\ $0.2548$ normalized). We therefore report the raw index
throughout, and an unqualified ``HHI'' means raw. Redemption outflow accompanied the de-peg: USDC's circulating supply fell from
\$43.29B on 7~March to \$35.90B on 18~March, a \emph{net} decline of \$7.39B, or $17.1\%$, over
the window. The \$8.15B above is \emph{gross} burns, which exceed the net decline because
minting continued through the window.

\begin{table}[t]
\centering
\caption{Redemption concentration, with an entity-resolved block. The top block reports
address-level figures, and the entity-resolved block joins the 776 funding addresses
against two public label sets (dawsbot/eth-labels, brianleect/etherscan-labels), which
together resolve 2 of 776 addresses: the larger conduit (40.9\% of volume) carries an
external Coinbase label, and the smaller (28.6\%) does not, so its Coinbase attribution
rests on the paper's own on-chain tracing. At this coverage no entity holds two or more of
the 776 addresses, so entity HHI equals address HHI to full precision, which means the
label evidence is too sparse to test the no-sybil condition.}
\label{tab:hhi}
\small
\checkwidth{\small% Source: results/number2_redemption_hhi.json --
%   upper block (unchanged from the current manuscript tab:hhi):
%     n_direct_senders; coverage.coverage_pct; upper_bound.{hhi_raw, hhi_normalized,
%     top10_share, largest_share, effective_n, bootstrap.{lo,hi}}; ex_custodian.{hhi_raw,
%     largest_share, effective_n}.
%   entity-resolved block (entity_resolved key):
%     labelled.union.{n_labelled, address_share} (coverage 2/776); externally_corroborated_share;
%     own_tracing_share; no_sybil_status; entity_hhi.{hhi_raw, bootstrap.{lo,hi}};
%     custodian_exclusion.all_labelled_exchange_custodian_excluded.hhi_raw;
%     custodian_exclusion.coinbase_conduits_excluded_existing_figure.hhi_raw (reproduces the
%     ex-custodian row above under entity resolution, for continuity -- same underlying figure,
%     not a second independent number); unlabelled_tail.{label_consistent_supremum_hhi,
%     label_consistent_supremum_assumption}.
% CONFIRMED: the old worst-case "0.577" figure does not appear anywhere in the live JSON
% (grep across usdc_depeg/results/ and paper/ returns no match) -- correctly
% identifies it as false, and it is absent, not merely uncited. Not reintroduced here.
% Width pass: Statistic/Value are both p{} (was p{8.4cm}r -- the old 'r' column
% could not wrap and let the 'No-sybil condition' prose row set the table's true
% width). Row labels trimmed to what the JSON itself supports; the longer
% constructions they drop (the entity-resolved section's own label-set names, the
% supremum row's 'entire unlabelled tail merged with...' clause) are unchanged in
% fc27.tex's existing tab:hhi caption, so no information is lost, only de-duplicated.
\begin{tabular}{@{}p{6.5cm}p{2.6cm}@{}}
\toprule
Statistic & Value \\
\midrule
Wallets (redeemers) & 776 \\
Coverage of gross burns & 89.8\% \\
HHI (raw) & 0.2558 \\
HHI (normalized) & 0.2548 \\
Top-10 share & 88.1\% \\
Largest share & 40.9\% \\
Effective \# redeemers ($1/\mathrm{HHI}$) & 3.9 \\
Bootstrap 95\% CI (HHI) & $[0.0517,\,0.4184]$ \\
\midrule
HHI, ex-custodian & 0.0671 \\
Largest share, ex-custodian & 16.8\% \\
Effective \# redeemers, ex-custodian & 14.9 \\
\midrule
\multicolumn{2}{@{}l}{\textit{Entity-resolved (public label join)}} \\
Labelled coverage & 2/776 \\
Externally corroborated share & 40.94\% \\
Own-tracing share & 28.63\% \\
Entity HHI (0 multi-address entities) & 0.2558 \\
No-sybil condition & untestable with public labels \\
HHI, labelled custodians excluded & 0.2528 \\
HHI, both Coinbase conduits excluded & 0.0671 \\
Supremum worst case & 0.9966 \\
\bottomrule
\end{tabular}
}
% Source: results/number2_redemption_hhi.json --
%   upper block (unchanged from the current manuscript tab:hhi):
%     n_direct_senders; coverage.coverage_pct; upper_bound.{hhi_raw, hhi_normalized,
%     top10_share, largest_share, effective_n, bootstrap.{lo,hi}}; ex_custodian.{hhi_raw,
%     largest_share, effective_n}.
%   entity-resolved block (entity_resolved key):
%     labelled.union.{n_labelled, address_share} (coverage 2/776); externally_corroborated_share;
%     own_tracing_share; no_sybil_status; entity_hhi.{hhi_raw, bootstrap.{lo,hi}};
%     custodian_exclusion.all_labelled_exchange_custodian_excluded.hhi_raw;
%     custodian_exclusion.coinbase_conduits_excluded_existing_figure.hhi_raw (reproduces the
%     ex-custodian row above under entity resolution, for continuity -- same underlying figure,
%     not a second independent number); unlabelled_tail.{label_consistent_supremum_hhi,
%     label_consistent_supremum_assumption}.
% CONFIRMED: the old worst-case "0.577" figure does not appear anywhere in the live JSON
% (grep across usdc_depeg/results/ and paper/ returns no match) -- correctly
% identifies it as false, and it is absent, not merely uncited. Not reintroduced here.
% Width pass: Statistic/Value are both p{} (was p{8.4cm}r -- the old 'r' column
% could not wrap and let the 'No-sybil condition' prose row set the table's true
% width). Row labels trimmed to what the JSON itself supports; the longer
% constructions they drop (the entity-resolved section's own label-set names, the
% supremum row's 'entire unlabelled tail merged with...' clause) are unchanged in
% fc27.tex's existing tab:hhi caption, so no information is lost, only de-duplicated.
\begin{tabular}{@{}p{6.5cm}p{2.6cm}@{}}
\toprule
Statistic & Value \\
\midrule
Wallets (redeemers) & 776 \\
Coverage of gross burns & 89.8\% \\
HHI (raw) & 0.2558 \\
HHI (normalized) & 0.2548 \\
Top-10 share & 88.1\% \\
Largest share & 40.9\% \\
Effective \# redeemers ($1/\mathrm{HHI}$) & 3.9 \\
Bootstrap 95\% CI (HHI) & $[0.0517,\,0.4184]$ \\
\midrule
HHI, ex-custodian & 0.0671 \\
Largest share, ex-custodian & 16.8\% \\
Effective \# redeemers, ex-custodian & 14.9 \\
\midrule
\multicolumn{2}{@{}l}{\textit{Entity-resolved (public label join)}} \\
Labelled coverage & 2/776 \\
Externally corroborated share & 40.94\% \\
Own-tracing share & 28.63\% \\
Entity HHI (0 multi-address entities) & 0.2558 \\
No-sybil condition & untestable with public labels \\
HHI, labelled custodians excluded & 0.2528 \\
HHI, both Coinbase conduits excluded & 0.0671 \\
Supremum worst case & 0.9966 \\
\bottomrule
\end{tabular}

\end{table}

\paragraph{Upper bound on ultimate-redeemer concentration.} The two largest funding addresses
($69.6\%$ of captured volume) are single-source routing conduits that are both funded by one
Coinbase hot wallet, which itself aggregates $14{,}111$ distinct senders. At the
presenting-entity level, Coinbase's $69.6\%$ implies an entity HHI near $0.49$. That
concentration measures a custodian's pass-through volume, but the claim here concerns
\emph{ultimate redeemers}. Disaggregating a custodian strictly lowers an HHI, so the
address-level figure is an upper bound on ultimate-redeemer concentration \emph{under a
no-sybil condition}: no entity outside the two dominant conduits splits its burns across
multiple funding addresses~\cite{Victor_2020}.

\begin{sloppypar}
Two public label sets (\texttt{dawsbot/eth-labels} and
\texttt{brianleect/etherscan-labels}) resolve only $2$ of the $776$ funding
addresses, so the no-sybil condition remains \emph{untestable}. At that coverage, the join
would detect a two-address entity with probability $3.33\times10^{-6}$. A null result from
that join is therefore uninformative. The larger conduit ($40.94\%$ of volume) carries the external label
``Coinbase: Circle Deposit 1''. The smaller ($28.63\%$) is unlabeled in both sets, and its
Coinbase attribution comes from our own on-chain tracing (Section~\ref{sec:setting}). The two sum to $69.57\%$, which the prose rounds to $69.6\%$
throughout. Excluding only the externally labeled custodian leaves the index at $0.2528$, and the
drop to $0.0671$ requires excluding the self-attributed, unlabeled conduit as well. We
therefore report $0.0671$ as conditional on our own attribution, and $0.2528$ as the
label-only floor. No informative worst-case cap follows from the label evidence: at that coverage, the
label-consistent supremum is $0.9966$. That supremum merges the entire unlabeled tail with the
one labeled custodian, so only our own attribution bounds the entity HHI.
\end{sloppypar}

Ma--Zeng--Zhang~\cite{mzz2025} report $521$ redeeming arbitrageurs per month for USDC
against $6$ for USDT, with USDC's largest single arbitrageur executing $45\%$ of monthly redemptions
against USDT's $66\%$. That is a different dataset and a different unit from our $40.9\%$ and
$69.6\%$ address-volume shares, so it does not replicate them. The
external label covers $40.9$ of the $69.6$ points of volume share. Within that share,
redemption ran through custodial conduits rather than through a few large arbitrageurs.
Beyond it, the same conclusion follows from our own attribution, with no external evidence
behind it.
